% =============================================================================
% Replacement for \subsection{Evidence Retrieval} (paper.tex, the old planner + beam-search section).
% Describes the LLM-free PPR retriever exactly as evaluated (nkg/retrieval/ppr_engine.py, setting t2000_f40_dmp06):
% no path completion, reserves, hub penalty, or second PPR round.
% New citations needed in references.bib: haveliwala2002topic, cormack2009reciprocal (entries at the end of this file).
% =============================================================================

\subsection{Evidence Retrieval}
\label{sec:retrieval}

At query time, \sys{} identifies the facts most relevant to a question $q$ by
propagating relevance over the index and then returns connected evidence
chains through those facts. Retrieval makes no LLM call. Following
HippoRAG~2~\cite{gutierrez2025rag}, it seeds Personalized PageRank
(PPR)~\cite{haveliwala2002topic} with nodes that match the question and lets
the graph spread their relevance to nodes that are connected to them but not
similar to the question. Unlike HippoRAG~2, whose propagation ranks
passages, \sys{} propagates over facts, entities, and chunks jointly and ranks
facts, which are then linked into paths over participation and semantic
edges. The reader receives these paths rather than isolated passages.

\paragraph{Seed selection.}
Let $\simcos$ denote cosine similarity between text embeddings, and let
$\bar{x}$ denote min--max normalization of a score over all nodes of one type.
\sys{} first retrieves the $K_c$ facts whose sentences are most similar to
$q$ and re-ranks them with a cross-encoder~\cite{nogueira2019passage}
$\mathrm{CE}(q,s_f)$, which scores the question and fact sentence jointly.
The top $N_s$ facts form the seed set $S$. The cross-encoder replaces the LLM
filtering of HippoRAG~2's recognition step: it removes facts that are similar to the
question without addressing it. Seeds then assign PPR reset mass to three node
types:
\begin{flalign*}
&u_f=\bar{c}_f\ \ (f\in S),\qquad
u_v=\frac{1}{|S_v|}\sum_{f\in S_v}\frac{\bar{c}_f}{|C_v|},\qquad
u_c=\beta\,\overline{\simcos}(q,t_c). &&
\end{flalign*}
Here, $c_f=\simcos(q,s_f)$ is a fact's sentence similarity,
$S_v=\{f\in S:(f,v)\in E_{FV}\}$ contains the seed facts in which entity $v$
participates, $C_v$ is the set of chunks containing facts about $v$, and
$t_c$ is the text of chunk $c$. All other facts and entities receive zero
reset mass, and the reset vector $u$ is normalized to sum to one. Seed facts contribute their
own relevance. Their participants inherit it, divided by the number of chunks in which
they occur, so specific entities seed strongly while generic entities, such
as a policy or an insured person mentioned throughout a contract, seed
almost nothing; HippoRAG~2 applies the same specificity weighting to
phrase nodes. Every chunk receives a small prior proportional to
its similarity to $q$, scaled by $\beta$ as in HippoRAG~2, so that
propagation can favour facts from relevant passages without letting chunks
dominate. In the policy example, the contestability rule $f_1$ and the
misstatement clause $f_4$ match the question, and \emph{date of birth},
which occurs in few chunks, inherits most of their entity mass.

\paragraph{Propagation.}
PPR runs on an undirected weighted graph over facts, entities, and chunks.
Participation edges $E_{FV}$ and provenance edges $E_{CF}$ have unit weight;
each semantic edge in $E_{FF}$ has weight equal to its extraction strength
$r$, taking the larger strength when both directions are present. With
$W$ the weighted adjacency matrix and $P$ its row-normalized transition
matrix, the stationary scores satisfy
\begin{flalign*}
&\pi=(1-d)\,u+d\,P^{\top}\pi, &&
\end{flalign*}
where $d$ is the probability of following an edge rather than restarting
at the seeds. Mass reaching a node without edges restarts from $u$, and $\pi$ is
obtained by power iteration. Chunks take part in propagation but never
appear in paths. Because entities and facts alternate, relevance reaches facts that share a
specific participant with a seed. Semantic edges also carry relevance to
facts that share no participant with a seed and are dissimilar to the
question, such as an exception or a condition stated in another
section. In the example, $f_2$ is reached through \texttt{has\_exception}
from $f_1$, and the age limit $f_5$ through the \texttt{depends\_on} edge
from $f_4$, although neither resembles the question.

\paragraph{Fact relevance.}
PPR ranks facts well in depth but can place facts that are merely central
near the top, whereas the cross-encoder is precise at the top but sees
only facts that resemble $q$. We combine the two by reciprocal-rank
fusion~\cite{cormack2009reciprocal}:
\begin{flalign*}
&\rho(f)=\frac{1}{k+\mathrm{rank}_{\pi}(f)}+\frac{1}{k+\mathrm{rank}_{\mathrm{CE}}(f)}. &&
\end{flalign*}
The cross-encoder scores a pool containing the $K_p$ facts with the highest
$\pi$ and the $K_p$ facts with the highest $c_f$. Facts outside the pool
receive only the PPR term. Rank fusion requires no calibration between
stationary probabilities and cross-encoder logits.

\paragraph{Path construction.}
The $K_f$ facts with the highest $\rho$ form the candidate set $F_q$.
Paths connect candidates through the index: a path extends from its last
fact to an unvisited fact in $F_q$ either directly over a semantic edge or
through a shared entity, $f_i\rightarrow v\rightarrow f_j$. When several
entities connect the same pair, the most specific one, with the fewest
facts, is used. A semantic edge takes precedence over an entity hop because
its label describes the connection. Restricting extensions to $F_q$ keeps
high-degree entities from multiplying the paths. A path scores
$\sigma(p)=\sum_{f\in p}\rho(f)$. A beam search starts from the top $N_0$
candidates, gives each partial path its $B$ highest-relevance extensions,
retains the $G$ best paths at each depth, and stops a path at $D$ facts or
when no unvisited candidate is reachable. Paths that reach $D$ facts, end
early, or fall outside the beam all remain candidates for selection.

\paragraph{Path selection.}
Many candidate paths repeat the same strong facts, so selection maximizes
the relevance of the evidence covered rather than the score of individual
paths. Given the set $\mathcal{C}$ of facts already covered, each step
selects the path with the largest gain
\begin{flalign*}
&\Delta(p\mid\mathcal{C})=\sum_{f\in p,\,f\notin\mathcal{C}}\rho(f), &&
\end{flalign*}
preferring shorter paths among ties, until $N_p$ paths are selected or no path
adds relevance. A short completion pass then ensures that each of the
$M$ most relevant facts appears: an uncovered fact contributes its
highest-gain candidate path, or the fact alone if it lies on no candidate
path. Coverage selection thus keeps chains that connect relevant facts
while still including relevant facts that the index leaves isolated.

\paragraph{Serialization.}
Selected paths are serialized in selection order. Facts are numbered and
quoted at their first mention and referred to by number thereafter; entities
appear as \texttt{[name (micro-role)]} at first mention and by name thereafter;
semantic hops carry their directional labels, and participation hops are
plain arrows. Steps already emitted by an earlier path are omitted, so each
path contributes only its new segments. The lowest-priority paths, with
completion paths last, are removed until the context fits a budget of $T$
tokens. Each fact appears once and remains traceable to its source chunks.
Optionally, the $P$ chunks with the highest PPR scores are appended
after the paths as source passages; this variant is marked separately in the
results. In the example, the highest-gain path
$f_1\xrightarrow{\texttt{has\_exception}}f_2\rightarrow\emph{date of birth}
\rightarrow f_4\xrightarrow{\texttt{depends\_on}}f_5$
covers the rule, its exception, the ineligible branch, and the age limit,
while a later path contributes only the alternative branch $f_3$.

\paragraph{Restricted views.}
The edge-restricted ablations apply the same restriction to propagation and
to paths. UNH-H removes semantic edges from $W$ and allows only entity hops;
UNH-D removes participation edges and allows only semantic hops. Provenance
edges and the chunk prior are retained in both views. In UNH-D, entities are
isolated, so mass placed on them restarts from the reset distribution.

Retrieval uses one embedding of the question, cross-encoder scoring of at
most $K_c+2K_p$ question--fact pairs, sparse power iteration, and a beam
search over at most $K_f$ facts. No retrieval step calls an LLM; answer generation
is a separate reader call.


% =============================================================================
% Replacement for the retrieval part of Table~\ref{tab:hyperparams} (keep the Construction block as is).
% =============================================================================
\midrule
\multicolumn{2}{@{}l}{\textit{Seed selection}} \\
Fact candidates re-ranked $K_c$ & 50 \\
Seed facts $N_s$ & 5 \\
Chunk prior weight $\beta$ & 0.05 \\
\midrule
\multicolumn{2}{@{}l}{\textit{Propagation and relevance}} \\
Edge-following probability $d$ & 0.6 \\
Participation, provenance, semantic edge weights & 1, 1, $r$ \\
Cross-encoder pool per ranking $K_p$ & 100 \\
Rank-fusion constant $k$ & 20 \\
\midrule
\multicolumn{2}{@{}l}{\textit{Path construction and selection}} \\
Candidate facts $K_f$ & 150 \\
Start facts $N_0$ & 30 \\
Extensions per partial path $B$ & 6 \\
Paths kept per depth $G$ & 60 \\
Maximum facts per path $D$ & 4 \\
Paths selected $N_p$ & 80 \\
Most relevant facts completed $M$ & 40 \\
Context budget $T$ (tokens) & 2,000 \\
Appended passages $P$ ($\ddagger$ rows only) & 5 \\

% Suggested replacement for the paragraph after Table~\ref{tab:hyperparams}:
% Table~\ref{tab:hyperparams} lists the \sys{} construction and retrieval settings, using the notation of
% Section~\ref{sec:method}. The same retrieval setting is used on every dataset. Sentence and chunk embeddings use
% voyage-4-nano, the cross-encoder is ms-marco-MiniLM-L6-v2, the budget $T$ is counted with the Qwen3 tokenizer, and
% power iteration runs for at most 60 iterations or until the $\ell_1$ change falls below $10^{-8}$. The reader is
% invoked separately after retrieval.

% =============================================================================
% New references.bib entries
% =============================================================================
% @inproceedings{haveliwala2002topic,
%   title={Topic-sensitive PageRank},
%   author={Haveliwala, Taher H.},
%   booktitle={Proceedings of the 11th International Conference on World Wide Web},
%   pages={517--526},
%   year={2002}
% }
%
% @inproceedings{cormack2009reciprocal,
%   title={Reciprocal rank fusion outperforms condorcet and individual rank learning methods},
%   author={Cormack, Gordon V. and Clarke, Charles L. A. and Buettcher, Stefan},
%   booktitle={Proceedings of the 32nd International ACM SIGIR Conference on Research and Development in Information Retrieval},
%   pages={758--759},
%   year={2009}
% }
